\section{通量、散度、旋度与边界定理}
\label{sec:02-Calculus/06-Multivariable-Calculus/09}

\subsection{水流穿过一张网：你能算，但不想每次从头来}

水以速度 \(\mathbf F=(x,0,0)\) 在空间中流动。你把一个单位立方体 \([0,1]^3\) 浸入水流中，想知道每秒有多少水从立方体里流出去。

直接算：左面 \(x=0\)，单位外法向量 \((-1,0,0)\)，\(\mathbf F\cdot\mathbf n = (0,0,0)\cdot(-1,0,0)=0\)，没有水穿过。右面 \(x=1\)，外法向 \((1,0,0)\)，\(\mathbf F\cdot\mathbf n = (1,0,0)\cdot(1,0,0)=1\)，通量 \(=1\times(\text{面积 }1)=1\)。另外四个面法向量与 \(\mathbf F\) 正交，通量为零。总通量：每秒 1 单位体积的水从右面流出。

六张平面，逐一做面法向量的点积，可行但枯燥。下一个问题是：把球面放进 \(\mathbf F=(x,y,z)\) 这个场里，通量是多少？球面上每个点的法向量都不同，每个点的 \(\mathbf F\) 也不同。你得对球面参数化、算叉积、做二重积分——一次性的力气能花，但换个曲面又要重来。

有没有一种方法，把穿过闭合曲面的总通量，转化成曲面内部某个量的累积，省掉在曲面上逐点积分的麻烦？

\subsection{通量：把场在曲面法向上的分量累加起来}

先建标准工具。设定向曲面由 \(\mathbf r(u,v)\) 参数化。曲面在每一点的法向面积微元是
\[
\mathbf r_u\times\mathbf r_v\,du\,dv.
\]
叉积的大小是局部面积，方向是曲面在该点的法向——哪个法向取决于参数 \((u,v)\) 的顺序。

向量场 \(\mathbf F\) 穿过这个曲面的通量为
\[
\iint_S\mathbf F\cdot\mathbf n\,dS
=\iint_D
\mathbf F(\mathbf r(u,v))
\cdot(\mathbf r_u\times\mathbf r_v)\,du\,dv.
\]
点积取出 \(\mathbf F\) 在法向上的分量——只有垂直于曲面的流动才``穿过''曲面，平行于曲面的流动在曲面上滑过，不贡献通量。交换参数顺序翻转法向和通量符号。闭曲面通常取外法向，但必须在问题里明确约定。

\subsection{散度：一个点的净流出率}

回到单位立方体和 \(\mathbf F=(x,0,0)\)。刚才算得净通量为 1。这个通量来自哪里？拿出来自右面的 1 和左面的 0——净通量完全来自 \(x\) 方向上 \(\mathbf F_x=x\) 随 \(x\) 增长而产生的不平衡。

把立方体缩小到一点：取包含 \(\mathbf p\) 的小体积 \(\Delta V\)，计算穿过其边界 \(\partial(\Delta V)\) 的通量，除以体积取极限：
\[
\nabla\cdot\mathbf F(\mathbf p)
=\lim_{\Delta V\to0}
\frac{1}{\Delta V}\iint_{\partial(\Delta V)}\mathbf F\cdot\mathbf n\,dS.
\]

对 \(\mathbf F=(P,Q,R)\)，直接计算给出
\[
\nabla\cdot\mathbf F=P_x+Q_y+R_z.
\]

散度是一个标量：正值在该点是源——净流出为正；负值在该点是汇——净流入为正。散度为零表示该点不可压缩——流入等于流出，没有净累积。

散度是两个截然不同概念的分界线：通量是穿过一块具体曲面的总量，散度是空间每一点的局部源强度。一块区域内部的散度累积，与穿过其边界面的总通量，应该能对上。

\subsection{Gauss 散度定理：内部源的积分等于穿过边界的总通量}

若 \(V\) 是适当光滑的体积，边界 \(\partial V\) 取外法向：
\[
\iiint_V\nabla\cdot\mathbf F\,dV
=\iint_{\partial V}\mathbf F\cdot\mathbf n\,dS.
\]

想一想这个等式如何出现。把 \(V\) 切成许多小块。相邻两块共享面上，一个取外法向 \(\mathbf n\)，另一个取 \(-\mathbf n\)，因此通量相反、互相抵消。所有内部面上的通量成对消失，只剩下外边界 \(\partial V\) 上的通量。而每小块的净流出又由其散度的体积分近似——小块足够多、足够细时，近似收敛到等式。

用开场的问题检验：\(\mathbf F=(x,y,z)\) 在半径为 \(R\) 的球内。散度：
\[
\nabla\cdot\mathbf F
=\frac{\partial x}{\partial x}+\frac{\partial y}{\partial y}+\frac{\partial z}{\partial z}
=1+1+1=3.
\]
散度处处为常值 3。体积分：
\[
\iiint_V3\,dV=3\cdot\frac43\pi R^3=4\pi R^3.
\]

球面上：\(\mathbf n=\mathbf r/R\)（单位径向），\(\mathbf F=(x,y,z)=R\mathbf n\)，所以 \(\mathbf F\cdot\mathbf n=R\)。通量：
\[
\iint_{\partial V}R\,dS
=R\cdot4\pi R^2=4\pi R^3.
\]
两端一致。散度定理把一次表面积分转化为了更简单的体积分——而且在这个例子里，转化后连积分都不必做，直接乘体积。

\subsection{旋度：局部旋转的向量表示}

散度测量源和汇，旋度测量``旋转''。对 \(\mathbf F=(P,Q,R)\)：
\[
\nabla\times\mathbf F = \begin{bmatrix}
R_y-Q_z\\
P_z-R_x\\
Q_x-P_y
\end{bmatrix}.
\]

旋度是一个向量。它的方向给局部旋转轴（按右手定则），大小与单位面积的环流正相关。

但``流线弯曲''不等于旋度非零。考虑 \(\mathbf F=(y,0,0)\)：流线是水平直线，但上层流得快、下层流得慢，上下层之间的剪切会产生旋度。再考虑 \(\mathbf F=(-y,x,0)/(x^2+y^2)\)（除原点外）：流线绕原点转圈，场看起来在旋转，但实际上在原点的任何不包含奇点的邻域内，旋度处处为零。旋度描述的是场在一点附近的反对称变化——不是单个粒子的轨迹曲率，不是流线的视觉弯曲程度，而是局部小回路上的环流密度。

\subsection{Stokes 定理：面上的旋转累积等于边界上的环流}

若定向曲面 \(S\) 的边界为 \(\partial S\)，方向按右手规则匹配：
\[
\iint_S(\nabla\times\mathbf F)\cdot\mathbf n\,dS
=\oint_{\partial S}\mathbf F\cdot d\mathbf r.
\]

验证的逻辑和散度定理同构：把曲面切成小片，内部公共边上的环流反向抵消，只剩外边界 \(\partial S\) 上的环流。每小片的局部环流由旋度的法向分量 \((\nabla\times\mathbf F)\cdot\mathbf n\) 近似。

平面上的特化就是 Green 定理：
\[
\oint_{\partial D}P\,dx+Q\,dy
=\iint_D(Q_x-P_y)\,dA,
\]
边界正向取``区域在行进方向左侧''。Green 定理也有通量形式，与二维散度 \(P_x+Q_y\) 对应：边界外法向通量的线积分，等于内部散度的面积分。

\subsection{三条定理是同一原则的不同维度}

从一维到三维，同一个``内部导数累积 = 边界净效应''的模式反复出现：

\begin{itemize}
\tightlist
\item 一维基本定理：\(\displaystyle\int_a^b f'(x)\,dx=f(b)-f(a)\)。导数在区间内累积，只剩端点差。
\item 线积分基本定理：\(\displaystyle\int_C\nabla\phi\cdot d\mathbf r=\phi(\text{终点})-\phi(\text{起点})\)。梯度沿路径积分，只剩端点势差。
\item Green / Stokes：旋度在面上累积，等于边界环流。
\item Gauss：散度在体内累积，等于边界通量。
\end{itemize}

统一表述是广义 Stokes 定理：
\[
\int_\Omega d\omega=\int_{\partial\Omega}\omega.
\]
外微分 \(d\) 在区域 \(\Omega\) 内部的累积，等于被积形式 \(\omega\) 在边界 \(\partial\Omega\) 上的积分。不同的 \(d\) 和 \(\omega\) 的组合，产生上述各个具体定理。

\subsection{两个恒等式：边界的边界是空的}

梯度场 \(\nabla\phi\) 取旋度：
\[
\nabla\times(\nabla\phi)=0.
\]
梯度场没有旋转——势场本身就是保守场。

旋度场 \(\nabla\times\mathbf F\) 取散度：
\[
\nabla\cdot(\nabla\times\mathbf F)=0.
\]
旋度场没有源或汇——旋转本身不产生净流出。

这两个恒等式有同一个几何根源。``边界的边界为空''：一条曲线的边界是两个端点，端点没有边界；一个曲面的边界是一条闭合曲线，闭合曲线没有端点。面到线到点，边界算子连续作用两次等于零。外微分形式中，\(d(d\omega)=0\) 精确地对应 \(\nabla\times(\nabla\phi)=0\) 和 \(\nabla\cdot(\nabla\times\mathbf F)=0\)。局部算子的连续作用后，下一层边界贡献互相抵消。

\subsection{条件、奇点与方向：定理不是免费的}

这三条定理要求场在包含了区域和边界的足够大的邻域内充分光滑。若区域内部有奇点，不能直接把定理套在包含奇点的区域上。正确的做法是：挖去奇点周围的小球或小柱，增加内边界，对剩余的``有洞''区域应用定理，再让挖去的部分趋于零——这一操作在后面的复分析和位势论中是标准套路。

方向错误只会差一个负号，但这个负号背后是一个未完整定义的几何对象：曲线行进方向、曲面法向和边界定向，必须成套选取。右手定则是标准约定，但换了左手规则只需整体一致地翻转符号——问题在于不配套的选择会悄无声息地给出正确大小但错误符号的结果。

延伸阅读参考 MIT OpenCourseWare 18.02SC Multivariable Calculus。

本篇的核心矢量分析工具在后续物理场论和流体力学建模中会反复出现；而``内部与边界''的思维方式，远超微积分，在凸分析、信息论和最优传输中都有对应结构。
